\vspace*{\stretch{1}}
\begin{center}
{\arabicfont\Large \beginR وَفَوْقَ كُلِّ ذِي عِلْمٍ عَلِيمٌ ۝٧٦\endR}

\vspace{2em}

{\large\itshape \enquote{And above every possessor of knowledge is \textbf{one} more knowing.}}

\vspace{1em}

{\small Quran, Surah Yusuf (12:76)}
\end{center}
\vspace*{\stretch{1.5}}
% \includepdf[pages=1,pagecommand={\thispagestyle{plain}}]{images/arabic_verse.pdf}
\cleardoublepage

\begin{abstract}
{\small
Software architecture documentation is often missing, incomplete, or outdated, making it unclear which source files implement which architectural modules. Software architecture recovery infers a modular decomposition from the implementation, while code-to-architecture mapping assigns files to known modules using a partial mapping. Both tasks can use dependencies, folder structure, and source code, but existing methods often rely on individual sources or predefined combinations. It remains unclear what each source contributes and when propagating information between connected files helps rather than mixes evidence from different modules.

The research uses open-source Java, C, and C++ systems with published reference decompositions. Recovery studies examine learning objectives, graph construction, and information sources, while mapping studies investigate representations, graph-based propagation, and iterative self-training.

For architecture recovery, learned representations achieve agreement with reference architectures comparable to or better than the evaluated baselines. In the evaluated configurations, reconstruction generally outperforms contrastive learning, while random walks and diffusion capture structural information beyond direct dependencies. However, combining structural, lexical, and folder information can improve agreement with reference architectures while producing modules with weaker dependency cohesion than those obtained from dependency-based random-walk embeddings.

For code-to-architecture mapping, neural classifiers outperform a Naive Bayes baseline. A graph attention network outperforms a graph-free classifier on most systems, although gains vary by system and configuration. Folder information produces the largest and most consistent improvements, but higher accuracy can coincide with misclassifications of files that a simple directory rule assigns correctly. Representations that encode dependencies generally outperform lexical ones, and dependency-guided fine-tuning improves contextual representations. 

However, strong graph-model performance does not by itself demonstrate a benefit from additional message passing. Controlled experiments show that propagating dependency information can reduce performance for some predictors compared with matched models without cross-file messages. Changes to graph construction, propagation methods, or learning settings can reduce these losses or lead to improvements.

Together, these findings show that the usefulness of architectural information depends not only on what is available, but also on how it is represented, combined, and used.

}

\bigskip
\vfill
\noindent
\textbf{Keywords:}
Software Architecture, Software Architecture Recovery, Representation Learning, Graph Neural Networks, Graph Autoencoders

\end{abstract}
\cleardoublepage


